% Generated draft: requires mathematical and source review.
\documentclass[11pt]{article}
\usepackage[margin=25mm]{geometry}
\usepackage{fontspec}
\setmainfont{Latin Modern Roman}
\usepackage{amsmath,amssymb,mathtools,unicode-math}
\setmathfont{Latin Modern Math}
\usepackage{enumitem,needspace,xurl,hyperref}
\setlength{\parindent}{0pt}
\setlength{\parskip}{5pt}
\setlength{\emergencystretch}{3em}
\newcommand{\N}{\mathbb{N}}
\newcommand{\Z}{\mathbb{Z}}
\newcommand{\Q}{\mathbb{Q}}
\newcommand{\R}{\mathbb{R}}
\newcommand{\C}{\mathbb{C}}
\newcommand{\F}{\mathbb{F}}
\newcommand{\A}{\mathbb{A}}
\newcommand{\PP}{\mathbb{P}}
\newcommand{\E}{\mathbb{E}}
\newcommand{\eps}{\varepsilon}
\newcommand{\dd}{\,\mathrm{d}}
\begin{document}
\section*{UnsolvedMath ID: 701405}
\textbf{Source problem code:} catalogue entry
\subsection{Definitions and mathematical statement}
Let \(\mathcal P\) denote the set of all primes. For \(\pi\subseteq\mathcal P\), a finite group is a \(\pi\)-group if every prime divisor of its order belongs to \(\pi\); the trivial group is included. A subgroup \(H\leq G\) of a finite group is a Hall \(\pi\)-subgroup if every prime divisor of \(|H|\) belongs to \(\pi\) and no prime divisor of \([G:H]\) belongs to \(\pi\).

Using the standard meaning of \(D_\pi\), a finite group \(G\) is a \(D_\pi\)-group if it has a Hall \(\pi\)-subgroup, any two Hall \(\pi\)-subgroups of \(G\) are conjugate in \(G\), and every \(\pi\)-subgroup of \(G\) is contained in a Hall \(\pi\)-subgroup of \(G\).

Define
\[
\mathcal S=\left\{\pi\subseteq\mathcal P:\ \text{for every finite group }G,\quad
\bigl(G\text{ has a Hall }\pi\text{-subgroup}\bigr)
\Longrightarrow
\bigl(G\text{ is a }D_\pi\text{-group}\bigr)
\right\}.
\]
The question proposed by D. O. Revin in 2010 is whether
\[
|\mathcal S|=2^{\aleph_0}.
\]
Here a continuum of sets means continuum many distinct subsets of \(\mathcal P\).

The supplied \emph{Comment of 2025} reports the following established special case, attributed to K. A. Ilenko and N. V. Maslova: for every integer \(x\in\Z\), the set
\[
\pi_x=\{p\in\mathcal P:p>x\}
\]
belongs to \(\mathcal S\). This reported result and the registry's open status are supplied source assessments, not independently verified here. The source does not explicitly define \(D_\pi\); the formulation above uses its standard group-theoretic definition and requires review against the cited notebook.
\subsection{Short English statement}
Are there continuum many sets \(\pi\) of primes such that every finite group with a Hall \(\pi\)-subgroup satisfies the full Sylow-type existence, conjugacy, and containment properties for \(\pi\)-subgroups? The supplied 2025 comment reports that every set of primes greater than an integer threshold has this property.
\subsection{Sources}
\begin{itemize}
\item Supplied local UnsolvedMath registry record: permanent ID \(701405\), exact source code \texttt{catalogue entry}, title \emph{Kourovka Notebook Problem 17.103}. The supplied status is \texttt{open}; the record was created and updated on \texttt{2026-10-16T14:00:00Z}. These are registry metadata, not an independent mathematical or literature assessment.
\item \emph{The Kourovka Notebook (Unsolved Problems in Group Theory)}, problem 17.103, proposed by D. O. Revin in 2010. The supplied background identifies the 17th issue (2010) and the 21st issue edition (2026), arXiv:1401.0300v48. The supplied collection description names E. I. Khukhro and V. D. Mazurov as editors. These source details have not been independently checked.
\item The supplied \emph{Comment of 2025} attributes the integer-threshold result to K. A. Ilenko and N. V. Maslova, \emph{Siberian Math. J.}, volume 62, no. 1 (2021), pp. 44--51. No article title or direct link is supplied; the result is retained as a source-reported qualification.
\end{itemize}
\end{document}
